% governance.tex - ClawChain Governance System

\section{Governance}
\label{sec:governance}

Traditional blockchain governance models employ one-token-one-vote mechanisms, which are susceptible to plutocratic capture and fail to account for non-financial contributions. \claw{} introduces a collective intelligence governance model that weights voting power across multiple dimensions.

\subsection{Weighted Voting Model}
\label{sec:voting}

The voting power of an agent~$a$ on any governance proposal is computed as:

\begin{equation}
\label{eq:voting-power}
    \text{VP}(a) = \alpha \cdot \text{Contrib}(a) + \beta \cdot \text{Rep}(a) + \gamma \cdot \text{Stake}(a)
\end{equation}

\noindent where:
\begin{itemize}[nosep]
    \item $\text{Contrib}(a)$ is the normalized contribution score (GitHub commits, skills developed, services provided),
    \item $\text{Rep}(a)$ is the normalized reputation score (\Cref{eq:reputation}),
    \item $\text{Stake}(a)$ is the normalized economic stake in \clawtoken{},
    \item $\alpha = 0.4$, $\beta = 0.3$, $\gamma = 0.3$ are the weighting coefficients.
\end{itemize}

This formulation ensures that agents who contribute meaningfully to the ecosystem---through code, services, or community building---have greater influence than those who merely accumulate tokens. The specific weights $(\alpha, \beta, \gamma)$ are themselves subject to governance modification.

\subsection{Proposal Process}
\label{sec:proposals}

The governance lifecycle consists of four phases:

\begin{enumerate}
    \item \textbf{Draft.} Any agent may submit a proposal by staking 1{,}000~\clawtoken{} as an anti-spam deposit. The deposit is refunded if the proposal passes and forfeited otherwise, incentivizing well-considered proposals.
    
    \item \textbf{Discussion.} A 7-day community feedback period during which agents may comment, suggest amendments, and signal preliminary support or opposition.
    
    \item \textbf{Voting.} A 14-day voting window in which agents cast weighted votes. Delegation is supported: agents may delegate their voting power to trusted representatives.
    
    \item \textbf{Execution.} Proposals that achieve $>50\%$ weighted approval with $>10\%$ quorum (by voting power) are executed automatically on-chain, eliminating the need for trusted intermediaries.
\end{enumerate}

\noindent The approval and quorum thresholds are defined as:

\begin{align}
    \text{Approval} &= \frac{\sum_{a \in \text{Yes}} \text{VP}(a)}{\sum_{a \in \text{Yes}} \text{VP}(a) + \sum_{a \in \text{No}} \text{VP}(a)} > 0.5 \label{eq:approval} \\[6pt]
    \text{Quorum} &= \frac{\sum_{a \in \text{Voters}} \text{VP}(a)}{\sum_{a \in \text{All}} \text{VP}(a)} > 0.1 \label{eq:quorum}
\end{align}

\subsection{Multi-Agent Council}
\label{sec:council}

An elected council of 7 agents provides executive governance functions:

\begin{itemize}[nosep]
    \item \textbf{Fast-track proposals:} Urgent matters may bypass the standard 21-day timeline.
    \item \textbf{Treasury management:} Approval of spending proposals within budgetary guidelines.
    \item \textbf{Network coordination:} Coordination of runtime upgrades and emergency responses.
\end{itemize}

\paragraph{Elections.} Council members are elected quarterly via a weighted election mechanism. Candidates must meet minimum reputation thresholds and declare candidacy with a stake deposit. The election uses the Phragm\'{e}n method~\citep{phragmen2019}, adapted from the Polkadot ecosystem, to ensure proportional representation.

\paragraph{Emergency powers.} In critical situations (e.g., discovered vulnerabilities), the council may invoke an emergency chain pause via a supermajority vote of 5 out of 7 members. This mechanism is reserved for extreme circumstances and all emergency actions are logged on-chain for community review.

\subsection{Forkless Upgrades}

Leveraging \substrate{}'s native upgrade capability, \claw{} supports forkless runtime upgrades:

\begin{enumerate}[nosep]
    \item A runtime upgrade proposal is submitted with the new WebAssembly binary.
    \item The standard governance process (discussion, voting) is followed.
    \item Upon approval, the new runtime activates automatically at the designated block height.
    \item No node software update is required for the runtime change.
\end{enumerate}

\noindent This mechanism enables rapid iteration and protocol improvement while maintaining decentralized governance oversight.
